\chapter*{M7: Shell}{}
 The goal of this milestone is to enable a user to interact with the system.
\addcontentsline{toc}{chapter}{M7: Shell}

\section{Command processing}

\subsection{Background tasks}
TODO: 
Have one main thread which periodically just spins in an event dispatch on a timer
(or another method) and checks serial in (which should be buffered to lines).
\\ 
If the command is a spawn and there is no \& make the main thread wait on the process(es)
\\
If there is an \& then create a new thread and have it wait on the process,
as an argument give it a number (n) and binary name. Have it print out [n] pid,
then wait on the process (the main thread should synchronize to not grab new input 
until this happens, we can use semaphores). then when it is done print [n] done and the command.
it may also make sense to write the output to a temp file that way we can print it all at once
when the process is done before we print done.
\\
this replicates bash behavior

\section{RamFS file system support}
We route all RamFS file operations through the \co{init} process. Each core has its own
in-memory file system. In order to support this we add file system functionality to
our RPC implementation. We add the following calls to the \co{aos_rpc} API:
\begin{lstlisting}
errval_t aos_rpc_fs_open(struct aos_rpc *chan, const char *path, int flags,
                         fs_fhandle_t *ret_handle);

errval_t aos_rpc_fs_close(struct aos_rpc *chan, fs_fhandle_t handle);

errval_t aos_rpc_fs_read(struct aos_rpc *chan, fs_fhandle_t handle,
                         void *buf, size_t nbytes, size_t *bytes_read);

errval_t aos_rpc_fs_write(struct aos_rpc *chan, fs_fhandle_t handle,
                          const void *buf, size_t nbytes, size_t *bytes_written);

errval_t aos_rpc_fs_seek(struct aos_rpc *chan, fs_fhandle_t handle,
                         enum fs_seekpos whence, off_t offset);

errval_t aos_rpc_fs_stat(struct aos_rpc *chan, fs_fhandle_t handle,
                         struct fs_fileinfo *info);

errval_t aos_rpc_fs_mkdir(struct aos_rpc *chan, const char *path);

errval_t aos_rpc_fs_rmdir(struct aos_rpc *chan, const char *path);

errval_t aos_rpc_fs_remove(struct aos_rpc *chan, const char *path);

errval_t aos_rpc_fs_opendir(struct aos_rpc *chan, const char *path,
                            fs_fhandle_t *ret_handle);

errval_t aos_rpc_fs_closedir(struct aos_rpc *chan, fs_fhandle_t handle);

errval_t aos_rpc_fs_readdir(struct aos_rpc *chan, fs_fhandle_t handle,
                            char **ret_name, struct fs_fileinfo *info);

\end{lstlisting}
The corresponding handlers are added to \co{rpc_server.c} and message tags to \co{aos_rpc.h}. The shell (or user process) can specify 
specify the flags which they want the file to be opened with, and can embedd the core which they
prefer to create the file on by passing \co{FS_OPEN_ON_CORE(core) | FS_OPEN_CREATE} to \co{aos_rpc_fs_open()} as the \co{flags} argument. This is intended to mirror
the spawn procedure of allowing a user to specify the core to spawn the procedure on. This works in the same way a 
shell spawn works (\cl{$ run process ---core n}), via the \cl{---core} command line argument like so,
\cl{$ touch file-name ---core n}. Any process can access files on any core, this is implemented with UMP redirection. To facilitate
demultiplexing file operations to the correct core, we embed the core in the \co{fs_handle_t} in the same manner
that the core of a process is embedded in its pid. Via this mechanism, \co{init} can route the operation to the correct
core without intervention from the user. One important constraint is that a user must supply the \co{FS_OPEN_ON_CORE(core)} flag to access 
cross core files (even when not creating them). This is because file systems on different cores do not share a namespace. This is 
an area for future work as the current implementation pushes some burden to the user.
\section{Pipes and I/O redirection}
To enable I/O redirection we change the spawn and RPC infrastructure. Processes are spawned
with information about where their input and output flows, then \co{aos_rpc_serial_getchar/putchar} correctly
routes characters (TODO: line granularity not characters) to the correct endpoint. Piping I/O between processes works 
in the following manner:
\begin{enumerate}
  \item The shell first creates and allocates a capability for the frame:
  \begin{lstlisting}
  struct capref pipe_frame;
  size_t pipe_size;
  frame_alloc(&pipe_frame, 2 * BASE_PAGE_SIZE, &pipe_size);
  \end{lstlisting}
  \item The shell then creates an \co{io_config_frame} for each side of the pipe and
  spawns the processes:
  \begin{lstlisting}
  // This omits the reader setup but it is nearly identical (the reader just sets pipe_in true and pipe_out_is_dsp 0)
  struct io_config_frame cfg = {
    .pipe_out        = true,
    .pipe_out_is_bsp = 1,
    .pipe_out_bytes  = pipe_size,
  };
  aos_rpc_proc_spawn_with_io(rpc, "writer_bin", core, pipe_frame, &cfg, &pid1);

  \end{lstlisting}
  \item Based on \co{core}, the correct \co{init} receives this \co{aos_rpc} call, and parses the configuration.
  \item \co{init} then creates a new frame for the IO config info and maps it into its virtual address space.
  \item \co{init} then assembles a capability array with capabilities to the pipe frame to pass to the spawned processes.
  \item \co{init} then spawns the process with \co{spawn_cmdline_with_caps()} which places the capabilities in the childs CSpace.
  \item In the new process, \co{barrelfish_init_onthread()} is called which does the necessary capability operations to map the pipe frame into its VSpace.
  \item This fills in a global variable identifying if I/O is piped or directed to a file.
  \item \co{aos_rpc_serial_getchar/putchar} correctly writes to the pipe rather than calling \newline \co{sys_putchar/getchar}.
\end{enumerate}
For file redirection the process is much the same except, fields are simply set in the \co{io_config_frame} that hold
\co{fs_handle_t} for the input and/or output file, no capability operations have to occur and are skipped with \co{NULL_CAP} checks. 
\section{Synchronizing serial input \& output}
TODO: Buffer calls to aos\_rpc\_serial\_getchar/putchar by only committing full lines (stop when newline char is seen) in aos\_rpc.c and route all calls to bsp init which can set
a global flag to serialize access. Make sure only bsp init can call sys\_putchar.
If this doesn't work find another way (maybe in the book?)
\section{Development process}
Shell development started when Shane implemented the infrastructure needed for file system 
operations over RPC, as well as the code which enables spawning processes with redirected I/O.